% ========================================================
%  《习近平新时代中国特色社会主义思想概论》  期末全真模拟试卷（第二套，含参考答案）
%  覆盖范围：导论、第一章至第十二章核心考点
%  题型：30单选(30分)+10多选(20分)+5简答(50分)
% ========================================================
\documentclass[12pt, a4paper]{article}

% ---- 中文支持 ----
\usepackage[UTF8]{ctex}

% ---- 数学 ----
\usepackage{amsmath, amssymb, amsthm}
\usepackage{mathtools}

% ---- 页面布局 ----
\usepackage[top=2.5cm, bottom=2.5cm, left=2.8cm, right=2.8cm]{geometry}

% ---- 表格 ----
\usepackage{booktabs, array, multirow, makecell}
\usepackage{tabularx}

% ---- 页眉页脚 ----
\usepackage{fancyhdr}
\usepackage{lastpage}
\pagestyle{fancy}
\fancyhf{}
\renewcommand{\headrulewidth}{0.6pt}
\lhead{\small 习近平新时代中国特色社会主义思想概论·期末全真模拟（第二套）}
\rhead{\small 共 \pageref{LastPage} 页}
\cfoot{\small 第 \thepage\ 页}

% ---- 计数器与样式 ----
\usepackage{enumitem}
\usepackage{xcolor}
\usepackage{tcolorbox}
\tcbuselibrary{skins, breakable}

% ---- 填空线 ----
\newcommand{\blank}[1]{\underline{\hspace{#1}}}

% ---- 答案盒子 ----
\newtcolorbox{answerbox}[1][]{
  breakable,
  colback=gray!6,
  colframe=gray!50,
  fonttitle=\bfseries\small,
  title=参考答案,
  #1
}

% ---- 题号格式 ----
\newcounter{bigq}
\newcommand{\BigQ}[1]{%
  \stepcounter{bigq}%
  \vspace{6pt}%
  \noindent{\large\bfseries 第\chinese{bigq}大题\quad #1}%
  \vspace{4pt}\par
}

\newcounter{subq}[bigq]
\newcounter{smallq}[subq]
\newcommand{\SmallQ}{%
  \stepcounter{smallq}%
  \noindent\arabic{smallq}.\quad
}
\newcommand{\resetsmallq}{\setcounter{smallq}{0}}

% 选项排版
\newcommand{\choices}[4]{\par
\noindent\hspace*{2em}A. #1\quad B. #2\quad C. #3\quad D. #4\par}

\newcommand{\choicelong}[4]{\par
\noindent\hspace*{2em}A. #1\par\hspace*{2em}B. #2\par\hspace*{2em}C. #3\par\hspace*{2em}D. #4\par}

\newcommand{\choicelongfive}[5]{\par
\noindent\hspace*{2em}A. #1\par\hspace*{2em}B. #2\par\hspace*{2em}C. #3\par\hspace*{2em}D. #4\par\hspace*{2em}E. #5\par}

% 分隔线
\newcommand{\examrule}{\noindent\rule{\linewidth}{0.6pt}}

% ---- 文档开始 ----
\begin{document}

% ============================================================
%  封面
% ============================================================
\begin{center}
  {\LARGE\bfseries 习近平新时代中国特色社会主义思想概论}\\[8pt]
  {\large\bfseries 期末全真模拟试卷（第二套）}\\[16pt]
  \begin{tabular}{lll}
    考试时间：120 分钟 & \quad & 满分：100 分 \\[4pt]
    \multicolumn{3}{l}{注意事项：请将答案写在答题纸指定区域，写在本卷上无效。}
  \end{tabular}
  \vspace{4pt}

  \examrule

  \vspace{6pt}
  \begin{tabular}{|c|c|c|c|c|}
    \hline
    \textbf{题号} & 一 & 二 & 三 & \textbf{总分} \\
    \hline
    \textbf{满分} & 30 & 20 & 50 & 100 \\
    \hline
    \textbf{得分} & & & & \\
    \hline
  \end{tabular}
\end{center}

\vspace{10pt}
\examrule
\vspace{6pt}

% ============================================================
%  第一大题：单项选择题（30题，每题1分，共30分）
% ============================================================
\BigQ{单项选择题（每题 1 分，共 30 分。每题只有一个正确选项）}

\resetsmallq
\SmallQ 习近平新时代中国特色社会主义思想创立的时代背景之一"中华民族伟大复兴进入关键时期"主要体现在（\quad）。
\choices{我国发展站在新的历史起点上，面临前所未有的机遇和挑战}{冷战格局加剧}{工业革命兴起}{经济全球化逆转}
\vspace{4pt}

\SmallQ 关于"两个确立"，下列表述正确的是（\quad）。
\choicelong{确立习近平同志党中央的核心、全党的核心地位；确立习近平新时代中国特色社会主义思想的指导地位}{确立党的基本路线；确立党的基本方略}{确立中国特色社会主义道路；确立理论自信}{确立"两个维护"；确立"两个结合"}
\vspace{4pt}

\SmallQ 中国特色社会主义进入新时代，这一论断回答的是我国发展新的（\quad）问题。
\choices{发展阶段}{历史方位}{历史阶段}{历史时期}
\vspace{4pt}

\SmallQ 中国共产党领导是中国特色社会主义最本质的特征，中国最大的国情是（\quad）。
\choices{中国共产党的领导}{人口规模巨大}{处于社会主义初级阶段}{发展中大国}
\vspace{4pt}

\SmallQ 在党的自身优势中，居于马克思主义政党第一位能力地位的是（\quad）。
\choices{思想引领力}{群众组织力}{政治领导力}{社会号召力}
\vspace{4pt}

\SmallQ 中国共产党的根本政治立场是（\quad）。
\choices{人民立场}{阶级立场}{国家立场}{民族立场}
\vspace{4pt}

\SmallQ 共同富裕是中国特色社会主义的（\quad），也是中国式现代化的重要特征。
\choices{本质要求}{基本制度}{现阶段任务}{长远规划}
\vspace{4pt}

\SmallQ 全面深化改革的总目标中，规定改革根本方向的是（\quad）。
\choicelong{建立社会主义市场经济体制}{实现全面建成小康社会}{构建新发展格局}{完善和发展中国特色社会主义制度}
\vspace{4pt}

\SmallQ 新发展理念中，注重解决发展不平衡问题的是（\quad）。
\choices{创新}{协调}{绿色}{开放}
\vspace{4pt}

\SmallQ "科技是第一生产力"的论断主要回答了（\quad）的关系问题。
\choices{科技与经济发展}{科技与国家强盛}{科技与文化交流}{科技与社会治理}
\vspace{4pt}

\SmallQ "人民民主是社会主义的生命"这一论断强调的是（\quad）。
\choices{民主的形态问题}{民主与社会主义的内在统一}{民主的程序问题}{民主的选举形式}
\vspace{4pt}

\SmallQ "两个结合"中，作为"根脉"的是（\quad）。
\choices{马克思主义基本原理}{中华优秀传统文化}{中国具体实际}{社会主义先进文化}
\vspace{4pt}

\SmallQ 习近平新时代中国特色社会主义思想科学体系中，体现世界观和方法论的部分是（\quad）。
\choices{十个明确}{十四个坚持}{十三个方面成就}{六个必须坚持}
\vspace{4pt}

\SmallQ 习近平新时代中国特色社会主义思想实现了马克思主义中国化时代化（\quad）。
\choices{新的飞跃}{第一次飞跃}{第二次飞跃}{第三次飞跃}
\vspace{4pt}

\SmallQ 道路问题是关系党的事业兴衰成败的第一位的问题，这一论断的核心在于（\quad）。
\choicelong{一个国家实行什么主义，关键要看能否解决其历史性课题}{道路决定社会性质}{道路决定经济体制}{道路决定文化形态}
\vspace{4pt}

\SmallQ "一个变与两个没有变"中，没有变的是（\quad）。
\choices{社会主要矛盾}{基本国情和国际地位}{社会性质和发展阶段}{经济体制和政治制度}
\vspace{4pt}

\SmallQ 党的基本路线可概括为（\quad）。
\choices{一个中心、两个基本点}{两个毫不动摇}{三个进一步解放}{四个全面}
\vspace{4pt}

\SmallQ 中国梦的内涵中，三个维度辩证统一，不包括（\quad）。
\choices{国家富强}{民族振兴}{人民幸福}{文化繁荣}
\vspace{4pt}

\SmallQ "五位一体"总体布局中，属于生态文明建设内容的是（\quad）。
\choicelong{推动形成绿色生产方式和生活方式}{发展全过程人民民主}{培育和践行社会主义核心价值观}{推进乡村振兴}
\vspace{4pt}

\SmallQ "五个必由之路"中，中国人民最显著的精神标识是（\quad）。
\choices{团结奋斗}{勇于创新}{改革开放}{自强不息}
\vspace{4pt}

\SmallQ 民心是最大的政治，因为民心（\quad）。
\choices{决定事业兴衰成败}{决定经济速度}{决定理论高度}{决定制度形态}
\vspace{4pt}

\SmallQ 贯彻群众路线的有效途径是（\quad）。
\choices{顶层设计}{调查研究}{试点推广}{集中决策}
\vspace{4pt}

\SmallQ 社会主义基本经济制度新概括中，作为所有制基础的是（\quad）。
\choices{两个毫不动摇}{按劳分配为主体}{社会主义市场经济体制}{公有制为主体、多种所有制经济共同发展}
\vspace{4pt}

\SmallQ 构建新发展格局最本质的特征是（\quad）。
\choices{以国内大循环为主体}{实现高水平的自立自强}{深化供给侧结构性改革}{扩大内需战略}
\vspace{4pt}

\SmallQ 全过程人民民主体现"四个相统一"，其中体现民主内容与形式统一的是（\quad）。
\choicelong{过程民主和成果民主的统一}{程序民主和实质民主的统一}{直接民主和间接民主的统一}{人民民主和国家意志的统一}
\vspace{4pt}

\SmallQ 伟大建党精神是（\quad）。
\choices{中国共产党人精神谱系的源头}{中华优秀传统文化的核心}{革命文化的全部}{社会主义先进文化的代表}
\vspace{4pt}

\SmallQ "让人民生活幸福是国之大者"体现了（\quad）。
\choices{以人民为中心的发展思想}{以经济建设为中心}{以改革为动力}{以稳定为前提}
\vspace{4pt}

\SmallQ 市域社会治理现代化的核心在于（\quad）。
\choicelong{市域是将风险矛盾化解在基层最直接最有效率的治理主体}{扩大城市规模}{增加财政投入}{提升行政级别}
\vspace{4pt}

\SmallQ 意识形态工作的主阵地、主战场、最前沿是（\quad）。
\choices{互联网}{传统媒体}{学校教育}{文艺创作}
\vspace{4pt}

\SmallQ "绿水青山就是金山银山"理念揭示了（\quad）的关系。
\choices{经济发展与生态环境保护}{工业与农业}{城市与乡村}{国内与国际}
\vspace{10pt}
\examrule

% ============================================================
%  第二大题：多项选择题（10题，每题2分，共20分）
% ============================================================
\BigQ{多项选择题（每题 2 分，共 20 分。每题有两个或两个以上正确选项，少选、错选均不得分）}

\resetsmallq
\SmallQ 习近平新时代中国特色社会主义思想的历史地位包括（\quad）。
\choicelong{当代中国马克思主义、二十一世纪马克思主义}{中华文化和中国精神的时代精华}{实现了马克思主义中国化时代化新的飞跃}{是党和国家必须长期坚持的指导思想}
\vspace{4pt}

\SmallQ 推进中国式现代化需要正确处理的重大关系包括（\quad）。
\choicelong{顶层设计与实践探索的关系}{战略与策略的关系}{守正与创新的关系}{效率与公平的关系}
\vspace{4pt}

\SmallQ 维护党中央权威和集中统一领导的重大意义包括（\quad）。
\choicelong{是党的领导的最高原则}{是马克思主义执政党成熟的重大建党原则}{与党的民主集中制完全一致}{意味着事事都要由党具体包办}
\vspace{4pt}

\SmallQ 中国式现代化的中国特色包括（\quad）。
\choicelong{人口规模巨大的现代化}{全体人民共同富裕的现代化}{物质文明和精神文明相协调的现代化}{人与自然和谐共生的现代化}
\vspace{4pt}

\SmallQ 新发展理念的实践要求包括（\quad）。
\choicelong{坚持创新在我国现代化建设全局中的核心地位}{在发展中促进相对平衡，不断增强发展的整体性}{推动形成绿色低碳的生产方式和生活方式}{推动形成更大范围、更宽领域、更深层次对外开放格局}
\vspace{4pt}

\SmallQ 健全党的全面领导制度的要求包括（\quad）。
\choicelong{完善党在各种组织中发挥领导作用的制度}{完善党协调各方的机制}{完善党领导各项事业的具体制度}{取消党对经济工作的领导}
\vspace{4pt}

\SmallQ 推动全体人民共同富裕取得更为明显的实质性进展，要求（\quad）。
\choicelong{共同富裕是中国特色社会主义的本质要求}{共同富裕是中国式现代化的重要特征}{坚持有底线、渐进式推进}{整齐划一、同步实现}
\vspace{4pt}

\SmallQ 全面深化改革开放的"三个进一步解放"是指（\quad）。
\choicelong{进一步解放思想}{进一步解放和发展社会生产力}{进一步解放和增强社会活力}{进一步解放私人资本}
\vspace{4pt}

\SmallQ 中国式现代化创造人类文明新形态的表现包括（\quad）。
\choicelong{深深植根于中华优秀传统文化}{体现科学社会主义的先进本质}{借鉴吸收一切人类优秀文明成果}{代表人类文明进步的发展方向}
\vspace{4pt}

\SmallQ 文化繁荣兴盛的重大意义在于（\quad）。
\choicelong{是实现中华民族伟大复兴的精神支撑}{是建设社会主义现代化强国的应有之义}{是满足人民日益增长的美好生活需要的内在要求}{是在世界文化激荡中站稳脚跟的前提基础}
\vspace{10pt}
\examrule

% ============================================================
%  第三大题：简答题（5题，每题10分，共50分）
% ============================================================
\BigQ{简答题（每题 10 分，共 50 分。要求观点正确、要点完整、简要论述）}

\resetsmallq
\SmallQ（10 分）简述习近平新时代中国特色社会主义思想是"两个结合"的重大成果。

\vspace{80pt}
\examrule
\vspace{6pt}

\SmallQ（10 分）简述"一个变与两个没有变"的内容及其辩证关系。

\vspace{80pt}
\examrule
\vspace{6pt}

\SmallQ（10 分）简述中国式现代化创造人类文明新形态的三个方面表现。

\vspace{80pt}
\examrule
\vspace{6pt}

\SmallQ（10 分）简述健全党的全面领导制度的具体要求。

\vspace{80pt}
\examrule
\vspace{6pt}

\SmallQ（10 分）简述新发展理念的科学内涵（五个方面对应解决的问题）。

\vspace{80pt}
\examrule

\vspace{16pt}
\examrule
\vspace{4pt}
\begin{center}
  {\large\bfseries ——试题到此结束，请认真检查——}
\end{center}
\vspace{4pt}
\examrule

% ============================================================
%  参考答案
% ============================================================
\newpage
\setcounter{bigq}{0}

\begin{center}
  {\LARGE\bfseries 参考答案与评分标准}
\end{center}
\vspace{8pt}
\examrule

% -------- 第一大题参考答案 --------
\BigQ{单项选择题参考答案}

\begin{answerbox}

\begin{tabular}{@{}llllll@{}}
1. A & 2. A & 3. B & 4. A & 5. C & 6. A \\
7. A & 8. D & 9. B & 10. B & 11. B & 12. B \\
13. D & 14. A & 15. A & 16. B & 17. A & 18. D \\
19. A & 20. A & 21. A & 22. B & 23. A & 24. B \\
25. B & 26. A & 27. A & 28. A & 29. A & 30. A \\
\end{tabular}

\vspace{6pt}
\textbf{要点解析：}

\begin{itemize}[leftmargin=2em,itemsep=2pt]
  \item 第1题：时代背景"五个方面"之一即中华民族伟大复兴进入关键时期，体现为我国发展站在新的历史起点上。
  \item 第2题："两个确立"标准表述——确立习近平同志党中央的核心、全党的核心地位；确立习近平新时代中国特色社会主义思想的指导地位。
  \item 第3题：新时代回答的是"历史方位"问题。
  \item 第4题：中国最大的国情就是中国共产党的领导。
  \item 第5题：政治领导力是马克思主义政党第一位的能力。
  \item 第6题：人民立场是中国共产党的根本政治立场。
  \item 第7题：共同富裕是中国特色社会主义的本质要求，也是中国式现代化的重要特征。
  \item 第8题：完善和发展中国特色社会主义制度规定改革根本方向；推进国家治理体系和治理能力现代化明确改革指向。题目问"规定根本方向"，故选D。
  \item 第9题：协调注重解决发展不平衡问题。
  \item 第10题："科技是第一生产力"主要回答科技与国家强盛的关系问题。
  \item 第11题："人民民主是社会主义的生命"强调民主与社会主义的内在统一。
  \item 第12题："根脉"是中华优秀传统文化，"魂脉"是马克思主义基本原理，注意区分。
  \item 第13题："六个必须坚持"是这一思想的世界观和方法论。
  \item 第14题：实现了马克思主义中国化时代化"新的飞跃"（注意：没有"第三次飞跃"的官方表述）。
  \item 第15题：道路问题核心在于"一个国家实行什么主义关键要看能否解决其历史性课题"。
  \item 第16题："两个没有变"是基本国情和国际地位。
  \item 第17题：党的基本路线可概括为"一个中心、两个基本点"。
  \item 第18题：中国梦三个维度是国家富强、民族振兴、人民幸福，不包括文化繁荣。
  \item 第19题：推动形成绿色生产方式和生活方式属于生态文明建设内容。
  \item 第20题：团结奋斗是中国共产党、中国人民最显著的精神标识。
  \item 第21题：民心是最大的政治，决定事业兴衰成败。
  \item 第22题：调查研究是贯彻群众路线的有效途径。
  \item 第23题："两个毫不动摇"是所有制基础。
  \item 第24题：构建新发展格局最本质的特征是实现高水平的自立自强。
  \item 第25题：程序民主和实质民主的统一体现民主内容与形式的统一。
  \item 第26题：伟大建党精神是中国共产党人精神谱系的源头。
  \item 第27题："让人民生活幸福是国之大者"体现以人民为中心的发展思想。
  \item 第28题：市域是将风险矛盾化解在基层、解决在萌芽状态最直接、最有效率的治理主体。
  \item 第29题：互联网是意识形态工作的主阵地、主战场、最前沿。
  \item 第30题："绿水青山就是金山银山"揭示了经济发展与生态环境保护的关系。
\end{itemize}

\end{answerbox}

% -------- 第二大题参考答案 --------
\BigQ{多项选择题参考答案}

\begin{answerbox}

\begin{tabular}{@{}ll@{}}
1. ABCD & 2. ABCD \\
3. ABC & 4. ABCD \\
5. ABCD & 6. ABC \\
7. ABC & 8. ABC \\
9. ABCD & 10. ABCD \\
\end{tabular}

\vspace{6pt}
\textbf{要点解析：}

\begin{itemize}[leftmargin=2em,itemsep=2pt]
  \item 第1题：思想历史地位四项表述全部正确。
  \item 第2题：推进中国式现代化需正确处理"六对重大关系"，本题列出四对全部正确。
  \item 第3题：维护党中央权威和集中统一领导是党的领导最高原则、是马克思主义执政党成熟的重大建党原则、与党的民主集中制完全一致；D项"事事由党具体包办"错误，党的领导是全面系统整体的，但不意味着事事包办。故选ABC。
  \item 第4题：中国式现代化五个中国特色，本题列出四个全部正确。
  \item 第5题：新发展理念五项实践要求全部正确。
  \item 第6题：健全党的全面领导制度三方面要求——完善党在各种组织中发挥领导作用的制度、完善党协调各方的机制、完善党领导各项事业的具体制度；D项"取消党对经济工作的领导"错误。故选ABC。
  \item 第7题：共同富裕是本质要求、重要特征，需坚持有底线、渐进式推进；D项"整齐划一、同步实现"错误。故选ABC。
  \item 第8题："三个进一步解放"指进一步解放思想、进一步解放和发展社会生产力、进一步解放和增强社会活力；D项错误。故选ABC。
  \item 第9题：人类文明新形态四方面表现全部正确。
  \item 第10题：文化繁荣兴盛四项意义全部正确。
\end{itemize}

\end{answerbox}

% -------- 第三大题参考答案 --------
\BigQ{简答题参考答案}

\begin{answerbox}

\textbf{第1题（10分）：习近平新时代中国特色社会主义思想是"两个结合"的重大成果}

习近平新时代中国特色社会主义思想是"两个结合"的重大成果，即坚持把马克思主义基本原理同中国具体实际相结合、同中华优秀传统文化相结合。（2分）

\textbf{（1）}坚持把马克思主义基本原理同中国具体实际相结合（3分）。这是"魂脉"所在。马克思主义基本原理必须同中国具体实际相结合，才能解决中国革命、建设、改革中的实际问题。习近平新时代中国特色社会主义思想立足新时代中国特色社会主义伟大实践，运用马克思主义立场、观点、方法回答了中国之问、世界之问、人民之问、时代之问。

\textbf{（2）}坚持把马克思主义基本原理同中华优秀传统文化相结合（3分）。这是"根脉"所在。中华优秀传统文化是中华民族的根和魂，是马克思主义在中国生根发芽的深厚文化土壤。新思想深刻汲取中华优秀传统文化的智慧结晶，把马克思主义思想精髓同中华优秀传统文化精华贯通起来，使马克思主义呈现出鲜明的中国特色、中国风格、中国气派。

\textbf{（3）}"两个结合"的统一（2分）。"魂脉"与"根脉"统一于马克思主义中国化时代化的实践，共同塑造了习近平新时代中国特色社会主义思想的鲜明理论品格。

\end{answerbox}

\begin{answerbox}

\textbf{第2题（10分）："一个变与两个没有变"的内容及其辩证关系}

\textbf{（1）"一个变"——社会主要矛盾的变化（3分）。}我国社会主要矛盾已经转化为人民日益增长的美好生活需要和不平衡不充分的发展之间的矛盾。这一变化是关系全局的历史性变化，对党和国家工作提出了许多新要求。

\textbf{（2）"两个没有变"（4分）。}
\begin{itemize}[leftmargin=2em,itemsep=1pt]
  \item 我国仍处于并将长期处于社会主义初级阶段的基本国情没有变。
  \item 我国是世界最大发展中国家的国际地位没有变。
\end{itemize}

\textbf{（3）辩证关系（3分）。}"一个变"与"两个没有变"是辩证统一的：社会主要矛盾的变化反映了我国发展阶段的局部性变化，但并没有改变我国社会主义所处历史阶段的判断；我国仍处于社会主义初级阶段这一最大实际没有变，决定了我们必须坚持党的基本路线不动摇；我国仍是世界最大发展中国家的国际地位没有变，决定了我国外交战略和对外方针的连续性。两者相结合，准确把握了我国发展所处的历史方位。

\end{answerbox}

\begin{answerbox}

\textbf{第3题（10分）：中国式现代化创造人类文明新形态的三个方面表现}

中国式现代化创造了人类文明新形态，主要体现在以下三个方面（每点约3分，共9分，逻辑结构1分）：

\textbf{（1）}提供了全新的现代化路径，超越西方现代化模式（3分）。中国式现代化打破了"现代化=西方化"的迷思，开辟了不同于西方现代化的新道路。它既坚持科学社会主义基本原则，又根据时代条件赋予其鲜明的中国特色，展现了现代化的多元可能性。

\textbf{（2）}为广大发展中国家提供了新的选择（3分）。中国式现代化成功证明，发展中国家完全可以立足自身国情探索现代化道路，不必照搬西方模式。这为希望既加快发展又保持自身独立性的国家和民族提供了全新选择，为人类文明进步贡献了中国智慧、中国方案。

\textbf{（3）}需牢牢把握推进中国式现代化的重大原则（3分）。创造人类文明新形态必须牢牢把握坚持和加强党的全面领导、坚持中国特色社会主义道路、坚持以人民为中心的发展思想、坚持深化改革开放、坚持发扬斗争精神这五大重大原则，确保中国式现代化行稳致远。

\end{answerbox}

\begin{answerbox}

\textbf{第4题（10分）：健全党的全面领导制度的具体要求}

坚持和加强党的全面领导，必须健全党的全面领导制度，发挥党在各种组织中总揽全局、协调各方的领导核心作用，把党的领导贯彻到党和国家所有机构履行职责全过程。具体要求包括以下三个方面（每点3分，共9分，总述1分）：

\textbf{（1）}必须完善党在各种组织中发挥领导作用的制度（3分）。要健全党领导人大、政府、政协、监察机关、审判机关、检察机关、武装力量、人民团体、企事业单位、基层群众性自治组织、社会组织等制度，确保党在各种组织中发挥领导核心作用。

\textbf{（2）}必须完善党协调各方的机制（3分）。要完善党统筹协调的体制机制，强化党的组织在同级组织中的领导地位，理顺党的组织同其他组织的关系，使党始终成为各项事业的坚强领导核心，不断增强党总揽全局、协调各方的能力。

\textbf{（3）}必须完善党领导各项事业的具体制度（3分）。要完善党领导经济社会发展各方面重要工作的制度，把党的领导贯彻到党和国家所有机构履行职责全过程，推动党领导各项事业的制度化、规范化、程序化，确保党中央重大决策部署落到实处。

\end{answerbox}

\begin{answerbox}

\textbf{第5题（10分）：新发展理念的科学内涵（五个方面对应解决的问题）}

新发展理念是创新、协调、绿色、开放、共享的发展理念，每项理念都对应解决我国发展中的特定问题（每点2分，共10分）：

\textbf{（1）}创新是引领发展的第一动力，\textbf{注重解决发展动力问题}。创新发展强调把创新摆在国家发展全局的核心位置，不断推进理论创新、制度创新、科技创新、文化创新等各方面创新。

\textbf{（2）}协调是持续健康发展的内在要求，\textbf{注重解决发展不平衡问题}。协调发展强调正确处理发展中的重大关系，推动区域协调发展、城乡协调发展、物质文明和精神文明协调发展。

\textbf{（3）}绿色是永续发展的必要条件和人民对美好生活追求的重要体现，\textbf{注重解决人与自然和谐共生问题}。绿色发展强调加快形成节约资源和保护环境的空间格局、产业结构、生产方式、生活方式。

\textbf{（4）}开放是国家繁荣发展的必由之路，\textbf{注重解决发展内外联动问题}。开放发展强调顺应我国经济深度融入世界经济的趋势，发展更高层次的开放型经济。

\textbf{（5）}共享是中国特色社会主义的本质要求，\textbf{注重解决社会公平正义问题}。共享发展强调坚持全民共享、全面共享、共建共享、渐进共享，不断推进全体人民共同富裕。

\end{answerbox}

\vspace{12pt}
\examrule
\begin{center}
  \textit{参考答案仅供参考，评分时应酌情给分，步骤正确可得过程分。}
\end{center}

\end{document}
